\documentclass[12pt]{article}
\usepackage[T1]{fontenc}
\usepackage[utf8]{inputenc}
\usepackage{lmodern}
\usepackage{textcomp}
\usepackage{enumitem}
\usepackage{geometry}
\geometry{margin=1in}
\begin{document}
\begin{center}
Answer Key - History - Graduate Level Assessment 5 \
\small (Answers are ordered exactly as in the question paper.)
\end{center}

\section*{Multiple Choice}
\begin{enumerate}[label=\arabic*.]
\item \textbf{Answer:} A
\\ \textit{Options:} A) Egypt; B) Western Europe; C) North America; D) South America; E) Eastern China
\\ \textit{Note:} The world's first civilization developed in Egypt due to the establishment of structured society, agriculture, and monumental architecture along the fertile banks of the Nile River, which fostered complex social, political, and economic systems around 3000 BCE.
\item \textbf{Answer:} A
\\ \textit{Options:} A) during the final minute of the film; B) 90 minutes into the film; C) about half way through the film; D) during the final 30 seconds of the film; E) in the first hour of the film
\\ \textit{Note:} The correct answer is based on the understanding that the history of the universe spans approximately 13.8 billion years, and the emergence of early human ancestors occurs only in the last moments of this extensive timeline, illustrating the relatively recent appearance of humans in the context of cosmic history.
\item \textbf{Answer:} C
\\ \textit{Options:} A) construction of permanent dwellings; B) invention of the wheel; C) parietal and mobiliary art; D) creation of writing systems; E) development of pottery
\\ \textit{Note:} The presence of parietal and mobiliary art during the Upper Paleolithic signifies a marked advancement in human intellect, as it reflects complex cognitive abilities, symbolic thinking, and the capacity for abstract representation, which are hallmarks of sophisticated cultural expressions.
\item \textbf{Answer:} D
\\ \textit{Options:} A) Canada.; B) Mexico.; C) the continental United States.; D) Alaska.
\\ \textit{Note:} Clovis points, distinctive stone tools associated with early Paleoindian cultures, have primarily been discovered in the continental United States but have not been found in Alaska, indicating a geographic limitation in their distribution.
\item \textbf{Answer:} D
\\ \textit{Options:} A) to demonstrate their military prowess to potential allies and enemies.; B) to expand their empire's territory.; C) to unite their own population through a common enemy.; D) to identify and cultivate talented leaders.; E) to gain control of vital resources.
\\ \textit{Note:} The Inca engaged in numerous wars not only for territorial expansion but also as a means to discover and develop capable military leaders, which was essential for maintaining and governing their vast empire effectively.
\end{enumerate}

\section*{True / False}
\begin{enumerate}[label=\arabic*.]
\setcounter{enumi}{5}
\item \textbf{Answer:} True
\\ \textit{Options:} A) True; B) False
\\ \textit{Note:} According to the passage: King Edward VII
\item \textbf{Answer:} False
\\ \textit{Options:} A) True; B) False
\\ \textit{Note:} The correct answer is: westward
\item \textbf{Answer:} False
\\ \textit{Options:} A) True; B) False
\\ \textit{Note:} The correct answer is: 1935/36
\item \textbf{Answer:} False
\\ \textit{Options:} A) True; B) False
\\ \textit{Note:} The correct answer is: video game handheld console market
\item \textbf{Answer:} True
\\ \textit{Options:} A) True; B) False
\\ \textit{Note:} According to the passage: Nero
\end{enumerate}

\section*{Long Form Response}
\begin{enumerate}[label=\arabic*.]
\setcounter{enumi}{10}
\item \textbf{Answer:} civil unrest and foreign invasions.
\\ \textit{Note:} Expected answer: civil unrest and foreign invasions.
\item \textbf{Answer:} Aerial Experiment Association
\\ \textit{Note:} Expected answer: Aerial Experiment Association
\end{enumerate}

\vfill
\noindent\textit{End of Answer Key}
\end{document}